% preambule.tex

% Русский язык через polyglossia
\usepackage{polyglossia}
\setmainlanguage{russian}
\newfontfamily{\cyrillicfonttt}{Courier New}

% Шрифты
\usepackage{fontspec}
\setmainfont{Times New Roman}

% Графика
\usepackage{graphicx}

% Фиксация float окружений
\usepackage{float}

% Математика
\usepackage{amsmath,amssymb,amsfonts,amsthm}

% Таблицы
\usepackage{makecell}
\usepackage{multirow}

% Листинги кода
\usepackage{listings}
\lstdefinestyle{mystyle}{
    basicstyle=\footnotesize\ttfamily,
    breaklines=true,
    numbers=left,
    numbersep=5pt,
    showstringspaces=false
}
\lstset{style=mystyle}

% Поля страницы
\usepackage{geometry}
\geometry{
    left=2.5cm,
    right=1.5cm,
    top=2cm,
    bottom=2cm
}

%%%%%%%%%%%%%%%%%%%%%%%%%%%%%%%%%%%%%%%%%%%%%%%% Всякое для ГОСТа

% Для подражания выравниванию по ширине 
\usepackage[none]{hyphenat}  % полное отключение переносов
\usepackage{microtype}       % улучшает выравнивание (важно!)
\emergencystretch=5em   % разрешаем растягивать строки
\tolerance=500          % терпимость к плохим строкам

% Межстрочный интервал
\usepackage{setspace}
\onehalfspacing 

% Отступы в списках
\usepackage{enumitem}
\setlist{
    itemsep=0pt,      % расстояние между пунктами
    parsep=0pt,       % расстояние между абзацами внутри пункта
    topsep=0pt,       % расстояние до списка сверху
    partopsep=0pt,    % дополнительное расстояние
    leftmargin=*,     % отступ текста
    labelindent=\parindent,  % отступ маркера списка
    align=left         % для labelindent
}

%%%%%%%%%%%%%%%%%%%%%%%%%%%%%%%%%%%%%%%%%%%%%%%% НАСТРОЙКА ЗАГОЛОВКОВ ПО МЕТОДИЧКЕ
\usepackage{titlesec}

% 1. Глава (\section) всегда начинается с новой страницы
\newcommand{\sectionbreak}{\clearpage}

% 2. Формат главы: 16пт, жирный, ПРОПИСНЫЕ
\titleformat{\section}
  {\fontsize{16}{20}\selectfont\bfseries\raggedright\MakeUppercase} 
  {\thesection} 
  {1em} 
  {} 

% 3. Интервалы главы: ДО - 0pt (т.к. новая страница), ПОСЛЕ - 24pt
\titlespacing*{\section}
  {0pt}   % отступ слева
  {0pt}   % интервал ДО
  {24pt}  % интервал ПОСЛЕ

% 4. Формат подраздела (\subsection): 14пт, жирный
\titleformat{\subsection}
  {\fontsize{14}{18}\selectfont\bfseries\raggedright}
  {\thesubsection}
  {1em}
  {}

% 5. Интервалы подраздела: ДО - 24pt, ПОСЛЕ - 12pt
\titlespacing*{\subsection}
  {0pt}   % отступ слева
  {24pt}  % интервал ДО
  {12pt}  % интервал ПОСЛЕ
%%%%%%%%%%%%%%%%%%%%%%%%%%%%%%%%%%%%%%%%%%%%%%%%

% Переопределяем название листингов
\renewcommand{\lstlistingname}{Листинг}

% Меняем формат caption
\usepackage{caption}
\captionsetup[lstlisting]{labelsep=endash} % тире для листингов

% Как форматировать листинги (язык по умолчанию)
\lstset{
    language=C++,
    frame=single,
    captionpos=b,            % подпись снизу            
    numbers=left,            % нумерация слева (внутри рамки)
    numbersep=5pt,           % расстояние от номера до кода
    xleftmargin=15pt,        % отступ слева (вместе с рамкой)
    framexleftmargin=15pt,   % отступ рамки от левого края
    xrightmargin=5pt,
    framexrightmargin=5pt
}

% Формат подписей рисунков (точка на тире)
\DeclareCaptionFormat{myformat}{#1 -- #3}
\captionsetup[figure]{name=Рисунок,format=myformat}

% Формат подписей таблиц (точка на тире)
\captionsetup[table]{labelsep=endash} 

% Для ссылок по ГОСТу
\usepackage{cite} 

% Теоремы, определения и прочее
\newtheorem{theorem}{Теорема}[section]
\newtheorem{lemma}{Лемма}[section]
\newtheorem{definition}{Определение}[section]
\newtheorem{proposition}{Утверждение}[section]
\newtheorem{corollary}{Следствие}[section]

% Углубленная структура секций (до 4 уровня)
\newcommand{\subsubsubsection}[1]{\paragraph{#1}\mbox{}\\}
\setcounter{secnumdepth}{4}
\setcounter{tocdepth}{4}

% Гиперссылки
\usepackage{hyperref}
\usepackage{bookmark}